\section{Introduction}

Financial markets generate multivariate data that are difficult to interpret using sector labels alone. Equity analysts commonly organise companies by sector classification, but firms within the same sector can differ in risk, liquidity, valuation, return behaviour and scale. Conversely, companies in different sectors may behave similarly because they share economic drivers, investor sentiment, balance-sheet structures or exposure to market cycles. These cross-sector relationships can be useful for portfolio screening, risk assessment and deeper analyst investigation.


Much of the machine learning literature in stock-market analysis focuses on supervised prediction, especially forecasting prices, returns or buy-sell signals. Such tasks are difficult because stock prices are noisy, non-stationary and influenced by macroeconomic news, firm announcements, liquidity changes and behavioural factors. A complementary question is not whether a model can predict the next price, but whether it can reveal hidden structure in a market at a point in time. This question is well suited to unsupervised learning because no outcome labels are required during model fitting.


This article studies the Johannesburg Stock Exchange (JSE) Top 40 as a South African emerging-market case. The JSE, established in 1887, is a major African equity exchange. The JSE Top 40 contains large, liquid companies across multiple sectors but remains small enough for detailed post-clustering interpretation. The objective is not to produce an automated trading rule. It is to evaluate whether a transparent unsupervised learning pipeline can identify coherent financial peer groupings and describe how those relationships differ across decision dates. This is important because several JSE Top 40 firms have global revenue, commodity, banking, telecommunications or retail exposure, making sector groupings useful but incomplete descriptions of multivariate financial similarity. In this repeated study, sector labels serve as an external descriptive comparator, not as verified historical ground truth.


The approach proposed by the paper is fourfold. First, the study extends the analysis to three decision dates, allowing financial peer groups to be compared over time. Second, UMAP is combined with hierarchical clustering, and cluster quality is evaluated in both the reduced representation and the scaled financial features. The earlier t-SNE hierarchical specification is retained as a comparator. Third, clusters are matched between dates using the membership of continuing instruments, so index entries and exits are separated from changes in peer-group membership. Fourth, cluster profiles, company-level feature changes, surrogate SHAP explanations and multivariate distance screening are combined to make the results easier to interpret. The research question is: Can interpretable unsupervised machine learning identify financial peer groups among JSE Top 40 equities, and how do these groups and their measured financial characteristics change across repeated cross-sections?


The article is positioned as exploratory financial data mining rather than predictive modelling. It does not claim to forecast future returns, determine fair value or generate a trading strategy. Its purpose is to show how interpretable unsupervised learning can support analyst attention by highlighting financial similarities, changes in peer relationships and unusual multivariate profiles. The repeated design also allows these findings to be examined in relation to index membership, feature definitions and model stability.


The remainder of the paper is structured as follows: Section 2 reviews related work, Section 3 describes the data and methodology, Section 4 presents the results, Section 5 discusses their implications, Section 6 outlines limitations and future work and Section 7 concludes.
